\begin{lemma}\label{lem:decic54-gonality}
The curve $X_1(54)$ has rational gonality at least $12$.
\end{lemma}
\begin{proof}
Put $X=X_1(54)$ and let $Y$ be the quotient by the order-$3$
projective diamond $\langle17\rangle$.
Then $g(X)=52$ and $g(Y)=10$.
For completeness, the congruence subgroup of $Y$ has index $324$,
no elliptic points, and $36$ cusps. Thus its genus is
$1+324/12-36/2=10$.
The cusp count can also be obtained directly: $X$ has $60$ geometric
cusps, of which $24$ are fixed by $\langle17\rangle$, so the quotient
has $(60+2\cdot24)/3=36$ cusps.
The exact primitive-coset calculation is recorded in
\path{parent_gonality_controls.json}.

The analytic bound \cite[(2.1)]{NV} gives
\[
 \operatorname{gon}_{\mathbb C}(X)
 \geq\frac{325\cdot972}{32768}>9.
\]
Suppose that $f:X\longrightarrow\mathbb P^1$ has degree
$n\in\{10,11\}$. Since $[\mathbb Q(X):\mathbb Q(Y)]=3$ is prime,
either $f$ factors through $Y$, or the two maps generate
$\mathbb Q(X)$. The first alternative would imply $3\mid n$.
In the second, Castelnuovo--Severi gives
\[
 52=g(X)\leq 3g(Y)+(3-1)(n-1)\leq50,
\]
a contradiction. This excludes both remaining degrees.
\end{proof}
